\chapter{Communication Complexity --- Entropy}
%%% generated by the blueprint pipeline from TCSlib.CommunicationComplexity.NewmanTheorem.Entropy
\section{Overview}

This module collects entropy and mutual-information lemmas needed for the Newman theorem
development.  It establishes upper bounds on (Shannon) entropy in terms of alphabet size,
monotonicity and data-processing inequalities for mutual information and conditional mutual
information, chain rules, and stability results under almost-sure equality and identical
distributions.

%%%%%%%%%%%%%%%%%%%%%%%%%%%%%%%%%%%%%%%%%%%%%%%%%%%%%%%%%%%%%%%%%%%%%%%%%%%%%%%
\section{Declarations}

\begin{theorem}[Entropy bounded by the logarithm of the alphabet size]
\lean{ProbabilityTheory.entropy_le_log_of_card_le}
\label{ProbabilityTheory.entropy_le_log_of_card_le}
\leanok
\difficulty{2}
Let $S$ be a finite nonempty set, let $(\Omega,\mu)$ be a measure space, and let
$X\colon\Omega\to S$ be a random variable. If $N$ is a natural number with $\abs{S}\le
N$, then the Shannon entropy of $X$ satisfies
\[
  H[X;\mu]\;\le\;\log N.
\]
\end{theorem}

\begin{theorem}[Entropy bound for an alphabet of at most $2^c$ symbols]
\lean{ProbabilityTheory.entropy_le_nat_mul_log_two_of_card_le_two_pow}
\label{ProbabilityTheory.entropy_le_nat_mul_log_two_of_card_le_two_pow}
\leanok
\uses{ProbabilityTheory.entropy_le_log_of_card_le}
\difficulty{1}
Let $S$ be a finite nonempty set and let $X$ be a random variable on a probability space
$\Omega$ (equipped with a measure $\mu$) taking values in $S$. If the number of elements
of $S$ is at most $2^c$ for some natural number $c$, then the Shannon entropy of $X$
satisfies
\[
  H[X;\mu] \;\le\; c\,\log 2,
\]
where the logarithm is natural (so the entropy is measured in nats).
\end{theorem}

\begin{theorem}[Mutual information is bounded by left entropy]
\lean{ProbabilityTheory.mutualInfo_le_entropy_left}
\label{ProbabilityTheory.mutualInfo_le_entropy_left}
\leanok
\difficulty{2}
Let $X$ and $Y$ be measurable random variables on a measure space $(\Omega,\mu)$, taking
values in countable measurable spaces and each attaining only finitely many values, and
suppose $\mu$ is either the zero measure or a probability measure. Then the mutual
information between $X$ and $Y$ is at most the Shannon entropy of $X$:
\[
  I[X : Y;\mu] \;\le\; H[X;\mu].
\]
\end{theorem}

\begin{theorem}[Mutual information is bounded by the entropy of either variable]
\lean{ProbabilityTheory.mutualInfo_le_entropy_right}
\label{ProbabilityTheory.mutualInfo_le_entropy_right}
\leanok
\uses{ProbabilityTheory.mutualInfo_le_entropy_left}
\difficulty{2}
Let $X \colon \Omega \to S$ and $Y \colon \Omega \to T$ be measurable random variables,
each taking only finitely many values, defined on a space $\Omega$ equipped with a
measure $\mu$ that is either the zero measure or a probability measure. Writing
$H[Y;\mu]$ for the Shannon entropy of $Y$ and $I[X:Y;\mu]$ for the mutual information
between $X$ and $Y$, one has
\[
  I[X:Y;\mu] \;\le\; H[Y;\mu].
\]
\end{theorem}

\begin{theorem}[Mutual information under almost-everywhere substitution]
\lean{ProbabilityTheory.mutualInfo_congr_ae}
\label{ProbabilityTheory.mutualInfo_congr_ae}
\leanok
\difficulty{1}
Let $(\Omega, \mu)$ be a measured space and let $S$ and $T$ be countable measurable
spaces. Suppose $X, X' \colon \Omega \to S$ and $Y, Y' \colon \Omega \to T$ are maps
with $X$ and $Y$ measurable, and that $X = X'$ and $Y = Y'$ hold $\mu$-almost
everywhere. Then the mutual information is unchanged, $I[X : Y \,;\, \mu] = I[X' : Y'
\,;\, \mu]$.
\end{theorem}

\begin{theorem}[Mutual information is invariant under injective recoding]
\lean{ProbabilityTheory.mutualInfo_comp_right_of_injective}
\label{ProbabilityTheory.mutualInfo_comp_right_of_injective}
\leanok
\difficulty{2}
Let $X\colon\Omega\to S$ and $Y\colon\Omega\to T$ be measurable random variables on a
probability space $(\Omega,\mu)$ (allowing also $\mu$ to be the zero measure), where $S$
and $T$ are countable and each of $X$ and $Y$ takes only finitely many values. If
$f\colon T\to V$ is a measurable injection into a countable space $V$, then passing $Y$
through $f$ leaves the mutual information with $X$ unchanged:
\[
I[X : f\circ Y\, ;\, \mu] = I[X : Y\, ;\, \mu].
\]
\end{theorem}

\begin{theorem}[Vanishing of mutual information for an almost-surely constant variable]
\lean{ProbabilityTheory.mutualInfo_eq_zero_of_ae_eq_const_left}
\label{ProbabilityTheory.mutualInfo_eq_zero_of_ae_eq_const_left}
\leanok
\difficulty{3}
Let $\mu$ be a probability measure (or the zero measure) on a sample space $\Omega$, and
let $X\colon\Omega\to S$ and $Y\colon\Omega\to T$ be measurable random variables taking
countably many values, each with finite range. If there is a constant $c\in S$ such that
$X=c$ holds $\mu$-almost everywhere, then the mutual information between $X$ and $Y$
vanishes: $I[X:Y;\mu]=0$.
\end{theorem}

\begin{theorem}[Vanishing of mutual information with an almost-surely constant variable]
\lean{ProbabilityTheory.mutualInfo_eq_zero_of_ae_eq_const_right}
\label{ProbabilityTheory.mutualInfo_eq_zero_of_ae_eq_const_right}
\leanok
\uses{ProbabilityTheory.mutualInfo_eq_zero_of_ae_eq_const_left}
\difficulty{2}
Let $X \colon \Omega \to S$ and $Y \colon \Omega \to T$ be measurable random variables
taking values in countable spaces, defined on a space $\Omega$ carrying a probability
measure $\mu$ (or the zero measure), and suppose each of $X$ and $Y$ takes only finitely
many values. If $Y$ is $\mu$-almost-surely equal to a fixed constant $c \in T$, then the
mutual information between $X$ and $Y$ vanishes: \[ I[X : Y \,;\, \mu] = 0. \]
\end{theorem}

\begin{theorem}[Vanishing conditional mutual information for an almost-surely constant variable]
\lean{ProbabilityTheory.condMutualInfo_eq_zero_of_ae_eq_const_left}
\label{ProbabilityTheory.condMutualInfo_eq_zero_of_ae_eq_const_left}
\leanok
\difficulty{3}
Let $\mu$ be a probability measure on a sample space $\Omega$ (or the zero measure), and
let $X$, $Y$, $Z$ be measurable random variables on $\Omega$ taking finitely many values
in the countable sets $S$, $T$, $U$ respectively. If there is a constant $c \in S$ with
$X = c$ holding $\mu$-almost everywhere, then the conditional mutual information of $X$
and $Y$ given $Z$ vanishes:
\[
I[X : Y \mid Z ; \mu] = 0.
\]
\end{theorem}

\begin{theorem}[Vanishing of conditional mutual information for an almost-surely constant variable]
\lean{ProbabilityTheory.condMutualInfo_eq_zero_of_ae_eq_const_right}
\label{ProbabilityTheory.condMutualInfo_eq_zero_of_ae_eq_const_right}
\leanok
\uses{ProbabilityTheory.condMutualInfo_eq_zero_of_ae_eq_const_left}
\difficulty{2}
Let $X$, $Y$, and $Z$ be discrete random variables on a probability space $(\Omega,
\mu)$ — where $\mu$ is either the zero measure or a probability measure — taking values
in countable measurable spaces $S$, $T$, and $U$ respectively, with each of $X$, $Y$,
$Z$ almost surely taking only finitely many values, and with $X$ and $Y$ measurable. If
there is a constant $c \in T$ such that $Y = c$ $\mu$-almost everywhere, then the
conditional mutual information satisfies $I[X : Y \mid Z; \mu] = 0$.
\end{theorem}

\begin{theorem}[Independence from factorization of fiber probabilities]
\lean{ProbabilityTheory.indepFun_of_measureReal_inter_preimage_singleton_eq_mul}
\label{ProbabilityTheory.indepFun_of_measureReal_inter_preimage_singleton_eq_mul}
\leanok
\difficulty{4}
Let $\mu$ be a finite measure on a measurable space $\Omega$, and let $S$ and $T$ be
finite sets. Let $X : \Omega \to S$ and $Y : \Omega \to T$ be measurable maps such that,
for every $x \in S$ and $y \in T$,
\[
\mu\bigl(X^{-1}\{x\} \cap Y^{-1}\{y\}\bigr) = \mu\bigl(X^{-1}\{x\}\bigr)\cdot
\mu\bigl(Y^{-1}\{y\}\bigr),
\]
where each measure is taken as a real number. Then $X$ and $Y$ are independent under
$\mu$.
\end{theorem}

\begin{theorem}[Conditional mutual information bounded by conditional entropy]
\lean{ProbabilityTheory.condMutualInfo_le_condEntropy_left}
\label{ProbabilityTheory.condMutualInfo_le_condEntropy_left}
\leanok
\difficulty{2}
Let $X$, $Y$, and $Z$ be random variables on a probability space $(\Omega, \mu)$, taking
values in countable measurable spaces, and assume each takes only finitely many values
(with positive probability). Then the conditional mutual information of $X$ and $Y$
given $Z$ is at most the conditional entropy of $X$ given $Z$:
\[
  I[X : Y \mid Z ; \mu] \;\le\; H[X \mid Z ; \mu].
\]
Here $\mu$ is allowed to be either a probability measure or the zero measure.
\end{theorem}

\begin{theorem}[Conditional mutual information bounded by conditional entropy]
\lean{ProbabilityTheory.condMutualInfo_le_condEntropy_right}
\label{ProbabilityTheory.condMutualInfo_le_condEntropy_right}
\leanok
\uses{ProbabilityTheory.condMutualInfo_le_condEntropy_left}
\difficulty{2}
Let $X$, $Y$, and $Z$ be measurable random variables on a common probability space
$(\Omega, \mu)$, each taking values in a countable measurable space and each having
finite range, and suppose $\mu$ is either the zero measure or a probability measure.
Then the conditional mutual information of $X$ and $Y$ given $Z$ is at most the
conditional entropy of $Y$ given $Z$:
\[
  I[X : Y \mid Z ; \mu] \;\le\; H[Y \mid Z ; \mu].
\]
\end{theorem}

\begin{theorem}[Conditional mutual information bounded by left entropy]
\lean{ProbabilityTheory.condMutualInfo_le_entropy_left}
\label{ProbabilityTheory.condMutualInfo_le_entropy_left}
\leanok
\uses{ProbabilityTheory.condMutualInfo_le_condEntropy_left}
\difficulty{1}
Let $X$, $Y$, and $Z$ be measurable random variables on a probability space $(\Omega,
\mu)$, taking values in countable spaces and each attaining only finitely many values.
Then the conditional mutual information of $X$ and $Y$ given $Z$ is bounded above by the
Shannon entropy of $X$:
\[
  I[X : Y \mid Z ; \mu] \;\le\; H[X ; \mu].
\]
(The bound also holds trivially when $\mu$ is the zero measure.)
\end{theorem}

\begin{theorem}[Conditional mutual information bounded by the entropy of one variable]
\lean{ProbabilityTheory.condMutualInfo_le_entropy_right}
\label{ProbabilityTheory.condMutualInfo_le_entropy_right}
\leanok
\uses{ProbabilityTheory.condMutualInfo_le_condEntropy_right}
\difficulty{1}
Let $X$, $Y$, and $Z$ be random variables on a probability space $(\Omega, \mu)$, taking
values in countable measurable spaces and each attaining only finitely many values. Then
the mutual information between $X$ and $Y$ conditioned on $Z$ is at most the Shannon
entropy of $Y$:
\[
  I[X : Y \mid Z; \mu] \;\le\; H[Y; \mu].
\]
\end{theorem}

\begin{theorem}[Conditional data processing with conditioning-dependent maps]
\lean{ProbabilityTheory.condMutualInfo_comp_left_le_of_comp_conditioning}
\label{ProbabilityTheory.condMutualInfo_comp_left_le_of_comp_conditioning}
\leanok
\difficulty{4}
Let $X$, $Y$, and $Z$ be random variables on a probability space $(\Omega, \mu)$, taking
values in countable measurable spaces $S$, $T$, and $U$ respectively, each measurable
and of finite range. For any function $f \colon U \times S \to V$ into a countable
measurable space $V$, form the postprocessed variable $W(\omega) = f\bigl(Z(\omega),
X(\omega)\bigr)$, obtained by applying to $X$ a map that may itself depend on the
conditioning value $Z$. Then the conditional mutual information satisfies
\[
  I[\,W : Y \mid Z;\, \mu\,] \;\le\; I[\,X : Y \mid Z;\, \mu\,].
\]
\end{theorem}

\begin{theorem}[Conditional data processing with conditioning-dependent maps]
\lean{ProbabilityTheory.condMutualInfo_comp_right_le_of_comp_conditioning}
\label{ProbabilityTheory.condMutualInfo_comp_right_le_of_comp_conditioning}
\leanok
\uses{ProbabilityTheory.condMutualInfo_comp_left_le_of_comp_conditioning}
\difficulty{3}
Let $X\colon\Omega\to S$, $Y\colon\Omega\to T$, and $Z\colon\Omega\to U$ be measurable
random variables on a probability space $(\Omega,\mu)$, each taking only countably many
values and having finite range, where $S$, $T$, and $U$ are countable measurable spaces.
Let $f\colon U\times T\to V$ be a measurable map into a countable measurable space $V$,
and define $W(\omega)=f\big(Z(\omega),Y(\omega)\big)$. Then the conditional mutual
information satisfies
\[
  I[X : W \mid Z;\mu] \;\le\; I[X : Y \mid Z;\mu].
\]
That is, replacing $Y$ by any function of $Y$ that is also permitted to depend on the
conditioning variable $Z$ cannot increase the conditional mutual information with $X$.
\end{theorem}

\begin{theorem}[Conditioning on $(Z, f(Z))$ versus conditioning on $Z$]
\lean{ProbabilityTheory.condMutualInfo_conditioning_prod_function_eq}
\label{ProbabilityTheory.condMutualInfo_conditioning_prod_function_eq}
\leanok
\difficulty{2}
Let $X$, $Y$, and $Z$ be discrete random variables on a probability space $(\Omega,
\mu)$, each taking values in a countable set and attaining only finitely many values,
and let $f$ be an arbitrary function defined on the range of $Z$. Then the conditional
mutual information of $X$ and $Y$ given the augmented variable $(Z, f(Z))$ equals the
conditional mutual information of $X$ and $Y$ given $Z$ alone:
\[
  I\bigl(X ; Y \mid (Z, f(Z))\bigr) = I(X ; Y \mid Z).
\]
\end{theorem}

\begin{theorem}[Adding a deterministic function of the pair as extra conditioning]
\lean{ProbabilityTheory.condMutualInfo_conditioning_prod_left_function_le}
\label{ProbabilityTheory.condMutualInfo_conditioning_prod_left_function_le}
\leanok
\difficulty{5}
Let $X$, $Y$, $Z$ be random variables on a probability space $(\Omega,\mu)$ (where $\mu$
is allowed to be the zero measure), taking values in countable measurable spaces $S$,
$T$, $U$ respectively, and assume each of $X$, $Y$, $Z$ is measurable with finite range.
Let $V$ be a further countable measurable space and let $f\colon U\times S\to V$ be
measurable, and write $f(Z,X)$ for the random variable $\omega\mapsto
f\bigl(Z(\omega),X(\omega)\bigr)$. Then conditioning additionally on $f(Z,X)$ cannot
increase the conditional mutual information of $X$ and $Y$ given $Z$:
\[
  I[X : Y \mid (Z,\, f(Z,X)) ; \mu]
  \;\le\;
  I[X : Y \mid Z ; \mu].
\]
\end{theorem}

\begin{theorem}[Conditioning on a function of the conditioned variables cannot raise conditional mutual information]
\lean{ProbabilityTheory.condMutualInfo_conditioning_prod_right_function_le}
\label{ProbabilityTheory.condMutualInfo_conditioning_prod_right_function_le}
\leanok
\uses{ProbabilityTheory.condMutualInfo_conditioning_prod_left_function_le}
\difficulty{2}
Let $X$, $Y$, $Z$ be random variables on a probability space $(\Omega, \mu)$, taking
values in countable measurable spaces, each measurable and with finite range. Let $f$ be
a measurable function of two arguments, and write $f(Z, Y)$ for the random variable
$\omega \mapsto f\bigl(Z(\omega), Y(\omega)\bigr)$. Then conditioning on the pair
$\bigl(Z, f(Z, Y)\bigr)$ — that is, refining the conditioning variable $Z$ by a
deterministic function of $Z$ and $Y$ — cannot increase the conditional mutual
information between $X$ and $Y$:
\[
  I\bigl[X : Y \mid (Z,\, f(Z, Y))\bigr] \;\le\; I\bigl[X : Y \mid Z\bigr].
\]
\end{theorem}

\begin{theorem}[Coarser conditioning does not increase conditional mutual information]
\lean{ProbabilityTheory.condMutualInfo_comp_conditioning_le_of_condMutualInfo_eq_zero}
\label{ProbabilityTheory.condMutualInfo_comp_conditioning_le_of_condMutualInfo_eq_zero}
\leanok
\difficulty{5}
Let $X$, $Y$, and $W$ be random variables on a probability space $(\Omega,\mu)$, taking
values in countable measurable spaces and each attaining only finitely many values, and
let $f$ be a measurable map defined on the range of $W$. Suppose the conditional mutual
information $I[X : W \mid f\circ W;\mu]$ vanishes, so that $X$ and $W$ are conditionally
independent given the coarsened variable $f\circ W$. Then
\[
  I[X : Y \mid f\circ W;\mu] \;\le\; I[X : Y \mid W;\mu].
\]
\end{theorem}

\begin{theorem}[Chain rule for conditional mutual information]
\lean{ProbabilityTheory.condMutualInfo_prod_left_eq_add}
\label{ProbabilityTheory.condMutualInfo_prod_left_eq_add}
\leanok
\difficulty{5}
Let $X$, $W$, $Y$, and $Z$ be random variables on a measure space $(\Omega,\mu)$, each
valued in a countable measurable space and each attaining only finitely many values, and
suppose $\mu$ is either the zero measure or a probability measure. Writing $I[U : V \mid
T;\mu]$ for the mutual information of $U$ and $V$ conditioned on $T$, the pair $(X,W)$
obeys the chain rule
\[
  I[(X,W) : Y \mid Z;\mu]
  \;=\; I[X : Y \mid Z;\mu] \;+\; I[W : Y \mid (X,Z);\mu].
\]
\end{theorem}

\begin{theorem}[Chain rule for conditional mutual information]
\lean{ProbabilityTheory.condMutualInfo_prod_right_eq_add}
\label{ProbabilityTheory.condMutualInfo_prod_right_eq_add}
\leanok
\uses{ProbabilityTheory.condMutualInfo_prod_left_eq_add}
\difficulty{2}
Let $X$, $Y$, $Z$, and $W$ be measurable functions on a probability space $(\Omega,
\mu)$, where $\mu$ is either the zero measure or a probability measure, each taking
values in a countable measurable space and attaining only finitely many values. Write
$I(X ; Y \mid Z)$ for the conditional mutual information between $X$ and $Y$ given $Z$.
Then the conditional mutual information between $X$ and the pair $(Y, W)$ given $Z$
satisfies
\[
I\bigl(X ; (Y, W) \mid Z\bigr) \;=\; I(X ; Y \mid Z) \;+\; I\bigl(X ; W \mid (Y,
Z)\bigr).
\]
\end{theorem}

\begin{theorem}[Adjoining a variable cannot decrease conditional mutual information]
\lean{ProbabilityTheory.condMutualInfo_le_prod_right_snd}
\label{ProbabilityTheory.condMutualInfo_le_prod_right_snd}
\leanok
\difficulty{2}
Let $X$, $Y$, $W$, and $Z$ be measurable functions on a probability space $(\Omega,
\mu)$, taking values in countable measurable spaces and each having finite range. Then
the conditional mutual information satisfies
\[
  I[X : W \mid Z; \mu] \;\le\; I[X : (Y, W) \mid Z; \mu],
\]
where $(Y, W)$ denotes the joint random variable $\omega \mapsto (Y(\omega),
W(\omega))$.
\end{theorem}

\begin{definition}[Strict prefix of a boolean vector]
\lean{ProbabilityTheory.boolVectorStrictPrefix}
\label{ProbabilityTheory.boolVectorStrictPrefix}
\leanok
For a boolean-vector-valued random variable $X : \Omega \to (\mathrm{Fin}\,m \to \mathrm{Bool})$
and an index $i : \mathrm{Fin}\,m$, the \emph{strict prefix} $\mathrm{boolVectorStrictPrefix}\,X\,i\,\omega$
is the restriction of $X(\omega)$ to coordinates $0, 1, \ldots, i-1$.
\end{definition}

\begin{theorem}[Chain rule for conditional mutual information over a boolean vector]
\lean{ProbabilityTheory.condMutualInfo_boolVector_eq_sum_strictPrefix}
\label{ProbabilityTheory.condMutualInfo_boolVector_eq_sum_strictPrefix}
\leanok
\uses{ProbabilityTheory.boolVectorStrictPrefix, ProbabilityTheory.condMutualInfo_eq_zero_of_ae_eq_const_left, ProbabilityTheory.condMutualInfo_prod_left_eq_add}
\difficulty{5}
Let $Y$ and $Z$ be random variables on a probability space $(\Omega, \mu)$ (allowing
also the zero measure), each taking countably many values with finite range, and let $X
: \Omega \to (\{0,1\}^m)$ be a random boolean vector; assume $X$, $Y$, and $Z$ are
measurable. Writing $X_i(\omega) = X(\omega)(i)$ for the $i$-th coordinate and
$X_{<i}(\omega)$ for the strict prefix consisting of coordinates $0, 1, \ldots, i-1$,
the conditional mutual information of $X$ and $Y$ given $Z$ decomposes as
\[
  I[X : Y \mid Z;\mu]
  \;=\; \sum_{i=0}^{m-1}
        I\bigl[X_i : Y \mid (X_{<i},\, Z);\mu\bigr].
\]
\end{theorem}

\begin{theorem}[Iterated conditioning on a subset]
\lean{ProbabilityTheory.cond_cond_eq_cond_of_subset}
\label{ProbabilityTheory.cond_cond_eq_cond_of_subset}
\leanok
\difficulty{2}
Let $\mu$ be a finite measure on a measurable space, and let $F \subseteq A$ be
measurable sets. Then conditioning $\mu$ on $A$ and afterwards on $F$ gives the same
measure as conditioning $\mu$ on $F$ directly:
\[
  (\mu[\,\cdot \mid A])[\,\cdot \mid F] \;=\; \mu[\,\cdot \mid F].
\]
\end{theorem}

\begin{theorem}[Multiplication rule for a nested conditioning event]
\lean{ProbabilityTheory.measureReal_mul_cond_real_eq_measureReal_of_subset}
\label{ProbabilityTheory.measureReal_mul_cond_real_eq_measureReal_of_subset}
\leanok
\difficulty{3}
Let $\mu$ be a finite measure on a measurable space $\Omega$, let $A$ be a measurable
set, and let $F \subseteq A$. Writing $\mu(\,\cdot \mid A)$ for the conditional measure
obtained by restricting $\mu$ to $A$ and normalising by $\mu(A)$, and taking all masses
as real numbers, one has
\[
  \mu(A)\,\cdot\,\mu(F \mid A) \;=\; \mu(F).
\]
\end{theorem}

\begin{theorem}[Chain decomposition of conditional mutual information over a finite variable]
\lean{ProbabilityTheory.condMutualInfo_prod_conditioning_eq_sum}
\label{ProbabilityTheory.condMutualInfo_prod_conditioning_eq_sum}
\leanok
\uses{ProbabilityTheory.measureReal_mul_cond_real_eq_measureReal_of_subset}
\difficulty{5}
Let $(\Omega,\mu)$ be a finite measure space, and let $X$, $Y$, $K$, and $Z$ be
measurable functions on $\Omega$ taking values in countable spaces, with $X$ and $Y$
each having finite range, $K$ taking values in a finite set $A$, and $Z$ taking values
in a finite space. Then the mutual information of $X$ and $Y$ conditioned on the pair
$(K,Z)$ decomposes as the $\mu$-weighted average, over the values of $K$, of the mutual
information of $X$ and $Y$ conditioned on $Z$ alone:
\[
  I[X : Y \mid (K,Z);\,\mu]
\;=\; \sum_{k \in A} \mu\bigl(K^{-1}\{k\}\bigr)\cdot I\bigl[X : Y \mid Z;\; \mu[\,\cdot
\mid K = k\,]\bigr],
\]
where $\mu(K^{-1}\{k\})$ denotes the (real-valued) measure of the event $\{K = k\}$ and
$\mu[\,\cdot \mid K = k\,]$ denotes $\mu$ conditioned on that event.
\end{theorem}

\begin{theorem}[Conditioning on a $Z$-event scales down conditional mutual information]
\lean{ProbabilityTheory.measureReal_mul_cond_condMutualInfo_le_condMutualInfo_of_event_eq_preimage}
\label{ProbabilityTheory.measureReal_mul_cond_condMutualInfo_le_condMutualInfo_of_event_eq_preimage}
\leanok
\uses{ProbabilityTheory.cond_cond_eq_cond_of_subset, ProbabilityTheory.measureReal_mul_cond_real_eq_measureReal_of_subset}
\difficulty{5}
Let $X$, $Y$, and $Z$ be random variables of finite range on a probability space
$(\Omega, \mu)$, taking values in countable measurable spaces, and let $A = Z^{-1}(B)$
be the event that $Z$ lands in a measurable set $B$. Writing $I[X : Y \mid Z;\, \nu]$
for the conditional mutual information of $X$ and $Y$ given $Z$ computed under a measure
$\nu$, and $\mu[\,\cdot \mid A]$ for $\mu$ conditioned on the event $A$, one has
\[
  \mu(A) \cdot I[X : Y \mid Z;\, \mu[\,\cdot \mid A]] \;\le\; I[X : Y \mid Z;\, \mu].
\]
\end{theorem}

\begin{theorem}[Chain rule for mutual information over a pair]
\lean{ProbabilityTheory.mutualInfo_prod_right_eq_add}
\label{ProbabilityTheory.mutualInfo_prod_right_eq_add}
\leanok
\difficulty{4}
Let $X$, $Y$, and $W$ be measurable random variables on a probability space $(\Omega,
\mu)$, taking values in countable measurable spaces, and suppose each of $X$, $Y$, $W$
takes only finitely many values almost surely. Then the mutual information between $X$
and the pair $(Y, W)$ decomposes as
\[
  I(X : (Y, W)) \;=\; I(X : Y) \;+\; I(X : W \mid Y),
\]
where $I(X : W \mid Y)$ is the mutual information of $X$ and $W$ conditioned on $Y$. (If
$\mu$ is the zero measure rather than a probability measure, both sides are $0$.)
\end{theorem}

\begin{theorem}[Conditioning monotonicity for mutual information]
\lean{ProbabilityTheory.condMutualInfo_le_mutualInfo_of_condDependence_le}
\label{ProbabilityTheory.condMutualInfo_le_mutualInfo_of_condDependence_le}
\leanok
\uses{ProbabilityTheory.mutualInfo_comp_right_of_injective, ProbabilityTheory.mutualInfo_prod_right_eq_add}
\difficulty{4}
Let $X$, $Y$, and $W$ be random variables on a probability space $(\Omega,\mu)$ (with
$\mu$ a probability measure, or else the zero measure), taking values in countable
measurable spaces and each assuming only finitely many values. Write $I[X : Y;\mu]$ for
the mutual information of $X$ and $Y$, and $I[X : Y \mid W;\mu]$ for their mutual
information conditioned on $W$. If conditioning on $Y$ does not increase the mutual
information between $X$ and $W$, that is,
\[
  I[X : W \mid Y;\mu] \;\le\; I[X : W;\mu],
\]
then conditioning on $W$ does not increase the mutual information between $X$ and $Y$:
\[
  I[X : Y \mid W;\mu] \;\le\; I[X : Y;\mu].
\]
\end{theorem}

\begin{theorem}[Conditional entropy is invariant under a.e.-equal random variables]
\lean{ProbabilityTheory.condEntropy_congr_ae}
\label{ProbabilityTheory.condEntropy_congr_ae}
\leanok
\difficulty{2}
Let $(\Omega,\mu)$ be a probability space, and let $X,X'\colon\Omega\to S$ and
$Y,Y'\colon\Omega\to T$ be measurable random variables valued in countable measurable
spaces, each taking only finitely many values. If $X = X'$ and $Y = Y'$ hold
$\mu$-almost everywhere, then the conditional Shannon entropies agree: \[ H[X \mid
Y;\mu] = H[X' \mid Y';\mu]. \]
\end{theorem}

\begin{theorem}[Conditional mutual information under a.e. change of both arguments]
\lean{ProbabilityTheory.condMutualInfo_congr_ae_left_right}
\label{ProbabilityTheory.condMutualInfo_congr_ae_left_right}
\leanok
\uses{ProbabilityTheory.condEntropy_congr_ae}
\difficulty{3}
Let $\mu$ be a probability measure on a sample space $\Omega$, and let $X, X' \colon
\Omega \to S$, $Y, Y' \colon \Omega \to T$, and $Z \colon \Omega \to U$ be measurable
maps into countable spaces, each taking only finitely many values $\mu$-almost surely.
If $X = X'$ and $Y = Y'$ hold $\mu$-almost everywhere, then the conditional mutual
informations agree:
\[
I[X : Y \mid Z\,;\,\mu] \;=\; I[X' : Y' \mid Z\,;\,\mu].
\]
\end{theorem}

\begin{theorem}[Conditional mutual information is invariant under a.e. equal variables]
\lean{ProbabilityTheory.condMutualInfo_congr_ae}
\label{ProbabilityTheory.condMutualInfo_congr_ae}
\leanok
\uses{ProbabilityTheory.condEntropy_congr_ae}
\difficulty{3}
Let $(\Omega,\mu)$ be a probability space, and let $X,X'$, $Y,Y'$, and $Z,Z'$ be
measurable random variables on $\Omega$ taking values in countable (discrete) measurable
spaces, each attaining only finitely many values. If $X = X'$, $Y = Y'$, and $Z = Z'$
hold $\mu$-almost everywhere, then the conditional mutual informations agree: \[ I[X : Y
\mid Z;\, \mu] = I[X' : Y' \mid Z';\, \mu]. \]
\end{theorem}

\begin{theorem}[Almost-everywhere invariance of conditional mutual information]
\lean{ProbabilityTheory.condMutualInfo_congr_ae_finite}
\label{ProbabilityTheory.condMutualInfo_congr_ae_finite}
\leanok
\uses{CommunicationComplexity.FiniteMeasureSpace, ProbabilityTheory.condMutualInfo_congr_ae}
\difficulty{2}
Let $\mu$ be a probability measure on a finite sample space $\Omega$ (with every subset
measurable), and let $X, X'$, $Y, Y'$, and $Z, Z'$ be random variables on $\Omega$
taking values in the finite sets $S$, $T$, and $U$ respectively. If $X = X'$, $Y = Y'$,
and $Z = Z'$ each hold $\mu$-almost everywhere, then the conditional mutual informations
coincide: \[ I[X : Y \mid Z\,;\,\mu] = I[X' : Y' \mid Z'\,;\,\mu]. \]
\end{theorem}

\begin{theorem}[Conditional mutual information depends only on the joint law]
\lean{ProbabilityTheory.IdentDistrib.condMutualInfo_eq}
\label{ProbabilityTheory.IdentDistrib.condMutualInfo_eq}
\leanok
\difficulty{4}
Let $X,Y,Z$ be measurable random variables on a probability space $(\Omega,\mu)$ and let
$X',Y',Z'$ be measurable random variables on a probability space $(\Omega',\mu')$, all
taking values in countable measurable spaces and each assuming only finitely many
values. If the triples $(X,Y,Z)$ and $(X',Y',Z')$ are identically distributed, meaning
they induce the same joint law on $S\times T\times U$, then their conditional mutual
informations agree: \[ I[X : Y \mid Z;\mu] = I[X' : Y' \mid Z';\mu']. \]
\end{theorem}

\begin{theorem}[Conditional mutual information depends only on the joint law]
\lean{ProbabilityTheory.IdentDistrib.condMutualInfo_eq_finite}
\label{ProbabilityTheory.IdentDistrib.condMutualInfo_eq_finite}
\leanok
\uses{CommunicationComplexity.FiniteMeasureSpace, ProbabilityTheory.IdentDistrib.condMutualInfo_eq}
\difficulty{2}
Let $(\Omega,\mu)$ and $(\Omega',\mu')$ be finite probability spaces, and let $S$, $T$,
$U$ be finite sets. Consider random variables $X\colon\Omega\to S$, $Y\colon\Omega\to
T$, $Z\colon\Omega\to U$ defined on the first space and $X'\colon\Omega'\to S$,
$Y'\colon\Omega'\to T$, $Z'\colon\Omega'\to U$ defined on the second. If the joint
random vectors $(X,Y,Z)$ and $(X',Y',Z')$ are identically distributed — that is, they
induce the same probability measure on $S\times T\times U$ — then their conditional
mutual informations agree: \[ I(X;Y\mid Z) = I(X';Y'\mid Z'). \]
\end{theorem}

\begin{theorem}[Conditional mutual information under injective recoding of the second and conditioning variables]
\lean{ProbabilityTheory.condMutualInfo_comp_right_conditioning_of_injective}
\label{ProbabilityTheory.condMutualInfo_comp_right_conditioning_of_injective}
\leanok
\difficulty{4}
Let $\mu$ be a probability measure on a sample space $\Omega$, and let $X$, $Y$, $Z$ be
measurable random variables on $\Omega$, taking values in countable measurable spaces
$S$, $T$, $U$ respectively and each having finite range. Let $f\colon T\to V$ and
$g\colon U\to W$ be measurable injections into countable measurable spaces. Then
relabelling the second and conditioning variables by $f$ and $g$ leaves the conditional
mutual information unchanged:
\[
  I[X : f\circ Y \mid g\circ Z ; \mu] \;=\; I[X : Y \mid Z ; \mu].
\]
\end{theorem}

\begin{theorem}[Identical distribution of a component under conditioning on the other]
\lean{ProbabilityTheory.IdentDistrib.cond_of_pair}
\label{ProbabilityTheory.IdentDistrib.cond_of_pair}
\leanok
\difficulty{5}
Let $X, Y$ be measurable maps from a probability space $(\Omega, \mu)$ taking values in
measurable spaces $A$ and $B$, and let $X', Y'$ be measurable maps from a probability
space $(\Omega', \mu')$ into the same spaces $A$ and $B$, with $Y$ and $Y'$ measurable.
Suppose the pairs $(X, Y)$ and $(X', Y')$ are identically distributed, that is, they
induce the same law on $A \times B$ under $\mu$ and $\mu'$ respectively. Then for every
measurable set $s \subseteq B$, the maps $X$ and $X'$ are identically distributed with
respect to the conditioned measures $\mu[\,\cdot \mid Y^{-1}(s)]$ and $\mu'[\,\cdot \mid
Y'^{-1}(s)]$.
\end{theorem}
